\section{Operator fusion 与 recomputation：少搬数据}

本节用工厂/仓库类比说明 HBM 往返为什么贵，再解释 fusion 和 recomputation 如何减少 memory access。


\begin{knowledgebox}{first-use glossary：operator fusion 与 recomputation}
Operator fusion 把多个小 kernel 合成一个 kernel，让中间结果留在寄存器/shared memory 中，减少 HBM 读写。Recomputation 是丢弃某些中间 activations，在需要时重算，用额外 compute 换更低 memory traffic 或 peak memory。
\end{knowledgebox}


这张 MXFP4 页也完成了第二部分到 operator fusion 的过渡：降低单元素 byte 只是减少 traffic 的一种方法，另一种方法是根本不把中间结果写回 HBM。于是下一节从“每次搬更少”转向“减少搬运次数”，比较 fusion 与 recomputation 如何改变 kernel 边界和 activation 生命周期。

\begin{knowledgebox}{讲义补充：源 Slide 29 应该怎么看}
这页是硬件机制或性能证据页。先看数据从哪里来、在哪里复用、是否写回 HBM；再判断瓶颈是 compute、memory bandwidth、thread scheduling 还是 alignment。GPU 优化不是让公式更漂亮，而是让数据在正确层级停留更久、让 tensor cores 少等数据。
\end{knowledgebox}


\begin{figure}[H]
\centering
\includegraphics[width=0.88\textwidth]{slide-030.jpg}
\caption{Slide 30: Trick 2: Operator fusion}
\figsource{CS336 Spring 2026 Lecture 5 官方幻灯片。}
\end{figure}

\noindent\textbf{展开说明：}Think of a GPU like a factory – inputs come from a warehouse (memory) and is

\begin{knowledgebox}{读图：Slide 30 应该怎么看}
这页是硬件机制或性能证据页。先看数据从哪里来、在哪里复用、是否写回 HBM；再判断瓶颈是 compute、memory bandwidth、thread scheduling 还是 alignment。GPU 优化不是让公式更漂亮，而是让数据在正确层级停留更久、让 tensor cores 少等数据。
\end{knowledgebox}


\begin{figure}[H]
\centering
\includegraphics[width=0.88\textwidth]{slide-031.jpg}
\caption{Slide 31: Operator fusion to minimize memory access}
\figsource{CS336 Spring 2026 Lecture 5 官方幻灯片。}
\end{figure}

\noindent\textbf{展开说明：}What if we have to do many operations? Shipping back and forth is somewhat silly

\begin{knowledgebox}{读图：Slide 31 应该怎么看}
这页是硬件机制或性能证据页。先看数据从哪里来、在哪里复用、是否写回 HBM；再判断瓶颈是 compute、memory bandwidth、thread scheduling 还是 alignment。GPU 优化不是让公式更漂亮，而是让数据在正确层级停留更久、让 tensor cores 少等数据。
\end{knowledgebox}


\begin{figure}[H]
\centering
\includegraphics[width=0.88\textwidth]{slide-032.jpg}
\caption{Slide 32: Example – sines and cosines}
\figsource{CS336 Spring 2026 Lecture 5 官方幻灯片。}
\end{figure}

\noindent\textbf{展开说明：}Computing sin 2 $x$ + cos 2 $x$ naively launches 5 CUDA kernels (back and forth)

\begin{knowledgebox}{读图：Slide 32 应该怎么看}
这页是硬件机制或性能证据页。先看数据从哪里来、在哪里复用、是否写回 HBM；再判断瓶颈是 compute、memory bandwidth、thread scheduling 还是 alignment。GPU 优化不是让公式更漂亮，而是让数据在正确层级停留更久、让 tensor cores 少等数据。
\end{knowledgebox}


\begin{figure}[H]
\centering
\includegraphics[width=0.88\textwidth]{slide-033.jpg}
\caption{Slide 33: Fusion example}
\figsource{CS336 Spring 2026 Lecture 5 官方幻灯片。}
\end{figure}

\noindent\textbf{展开说明：}All 5 pointwise operations can be fused into a single CUDA kernel call.


\begin{figure}[H]
\centering
\includegraphics[width=0.88\textwidth]{slide-034.jpg}
\caption{Slide 34: Trick 3: recomputation}
\figsource{CS336 Spring 2026 Lecture 5 官方幻灯片。}
\end{figure}

\noindent\textbf{展开说明：}[From cs221]

\begin{knowledgebox}{读图：Slide 34 应该怎么看}
这页是硬件机制或性能证据页。先看数据从哪里来、在哪里复用、是否写回 HBM；再判断瓶颈是 compute、memory bandwidth、thread scheduling 还是 alignment。GPU 优化不是让公式更漂亮，而是让数据在正确层级停留更久、让 tensor cores 少等数据。
\end{knowledgebox}


\begin{figure}[H]
\centering
\includegraphics[width=0.88\textwidth]{slide-035.jpg}
\caption{Slide 35: Storing (and retrieving) activations can be expensive!}
\figsource{CS336 Spring 2026 Lecture 5 官方幻灯片。}
\end{figure}

\noindent\textbf{展开说明：}Let’s say we stack 3 sigmoids on top of each other.

\begin{knowledgebox}{读图：Slide 35 应该怎么看}
这页是硬件机制或性能证据页。先看数据从哪里来、在哪里复用、是否写回 HBM；再判断瓶颈是 compute、memory bandwidth、thread scheduling 还是 alignment。GPU 优化不是让公式更漂亮，而是让数据在正确层级停留更久、让 tensor cores 少等数据。
\end{knowledgebox}

\subsection{本章小结}
本节的共同问题是如何减少昂贵的数据移动并提高硬件利用率。GPU 的速度来自海量并行和专用矩阵硬件，但只有当 memory access、shape、precision 和 kernel 组织匹配时，这些速度才会真正出现。

